\originalchapter{第一次期中考试}
\psource{8.044，2013 春，Exam 1（官方档名 e1\_b），三题 30/30/40 分。中文题面；据公开官方答案核对的 AI 独立解答。}
\originalsection{1．量子点（30 分）}
\begin{exercise}
边长 $L$ 的硅片上量子点密度 $p(x,y)=(x+y)/L^3$，$0\le x,y\le L$，其他处零。
(a，15 分) 求 $p(x\mid y)$，画 $y=L/2$ 情况，判断独立性并解释；(b，15 分) 求 $M=\max(x,y)$ 密度并画图，无须前问。
\end{exercise}
\begin{solution}
(a) $p_y=(y+L/2)/L^2$，$p(x\mid y)=(x+y)/[L(y+L/2)]$。$y=L/2$ 时为 $(x+L/2)/L^2$，端点 $1/(2L),3/(2L)$。条件密度依赖 $y$，不独立。
(b) $F_M(m)=\int_0^m\int_0^m(x+y)\dd x\dd y/L^3=(m/L)^3$，$p_M=3m^2/L^3$（$0\le m\le L$）。端外密度零。
\end{solution}
\begin{center}\begin{tikzpicture}\begin{axis}[width=6.5cm,height=3.6cm,xmin=0,xmax=1,ymin=0,ymax=3.1,xlabel={归一化位置},ylabel={$L$ 倍密度},legend style={font=\small}]\addplot[domain=0:1]{x+.5};\addlegendentry{条件密度}\addplot[domain=0:1,thick]{3*x*x};\addlegendentry{最大值密度}\end{axis}\end{tikzpicture}\end{center}
\originalsection{2．真实气体绝热线（30 分）}
\begin{exercise}
模型 $U=aV^{-2/3}+bV^{2/3}T^2$、$P=2aV^{-5/3}/3+2bV^{-1/3}T^2/3$，$a,b$ 包含固定总原子数的依赖。求通过 $(T_0,V_0)$ 的绝热路径 $V(T)$。
\end{exercise}
\begin{solution}
准静态无热 $\dd U+P\dd V=0$，$a$ 项相消，余 $2bV^{2/3}T\dd T+4bV^{-1/3}T^2\dd V/3=0$。对 $b,T>0$，$\dd T/T=-2\dd V/(3V)$，所以 $V=V_0(T_0/T)^{3/2}$。不可逆绝热过程还需额外过程信息，不能直接指定此曲线。
\end{solution}
\originalsection{3．一维超相对论气体（40 分）}
\begin{exercise}
$N$ 个相同经典微弱耦合原子，$0\le x\le L$、动量大小 $p\ge0$、方向 $s=\pm1$，$\epsilon=cp$。微正则总能 $E$ 至 $E+\Delta$，$\Delta\ll E$。
(a，16 分) 求无量纲可达相空间，按题给以 $\hbar$（长度乘动量）作态单元，并防 Gibbs 悖论；可用正单纯形体积 $S^d/d!$；(b，8 分) 求 $E(T,L,N)$，熵中 $\ln(N\Delta/E)$ 可略；(c，8 分) 求张力 $\mathcal F$，其负值相当于压力向外；(d，8 分) 求过 $(T_0,L_0)$ 的绝热线 $L(T)$。
\end{exercise}
\begin{solution}
(a) 两方向各粒子给 $2^N$，位置 $L^N$，大小动量为正单纯形，且相同粒子除 $N!$：
\[\Phi=\frac1{N!}\left(\frac{2L}{\hbar c}\right)^N\frac{E^N}{N!},\quad
\Omega\simeq\frac1{N!}\left(\frac{2L}{\hbar c}\right)^N\frac{E^{N-1}\Delta}{(N-1)!}.\]
两阶乘来源不同；标准量子态单元若换 $h$ 仅改熵常数。
(b) 领先 $S\simeq N\kb\ln[2L E e^2/(\hbar cN^2)]$，故 $E=N\kb T$；固定宽壳精确导数给 $(N-1)\kb T$。
(c) $\dd E=T\dd S+\mathcal F\dd L$，$\mathcal F=-T S_L=-N\kb T/L$。
(d) 可逆等熵要求 $LE$ 不变（$N$ 固定），$LT=L_0T_0$，即 $L=L_0T_0/T$。
\end{solution}
